\section{Implementation}\label{sec:implementation}

This chapter describes the implemented proof-of-concept system. The implementation follows the narrowed Tier 1.5 scope defined for the thesis: one temperature-event path, two workflow decision modes, one shared action contract, one minimum safety and validation design, simulated fan actions, CSV-based evaluation, and Raspberry Pi validation limited to the middleware/action endpoint side. The purpose of the implementation was not to build a broad automation platform, a multi-agent architecture, or a production safety framework. Instead, the implementation creates a complete and measurable vertical slice that can be evaluated against the research questions.

\subsection{System Overview}\label{subsec:implementation-system-overview}

The final frozen system is organized around a temperature event that enters a workflow and is represented as a contract-shaped fan action object before reaching the middleware. At the highest level, the implemented pipeline is:

\begin{quote}
temperature event/test input \(\rightarrow\) n8n workflow \(\rightarrow\) workflow decision \(\rightarrow\) shared contract and documented validation boundary \(\rightarrow\) Python middleware \(\rightarrow\) simulated \texttt{fan\_on} or \texttt{fan\_off} action \(\rightarrow\) evaluation result files.
\end{quote}

The architecture deliberately separates workflow logic from action execution. n8n is responsible for receiving the event and selecting an action path. The shared JSON contract describes the message shape that workflow and agent output must follow. The safety-layer files define the minimum policy for valid, malformed, and risky action output. The Python middleware then acts as the controlled API boundary for action execution. In the frozen implementation, the action endpoint changes an in-memory simulated fan state rather than controlling physical hardware.

Figure~\ref{fig:implementation-final-architecture} shows the final report architecture. The figure also shows the Tier 1.5 Raspberry Pi validation path. In that validation, n8n remained on the PC, while only the Python middleware/action endpoint ran on the Raspberry Pi. The PC-hosted n8n workflow called the Raspberry Pi middleware over HTTP. This preserves the PC-first architecture while still testing whether the action endpoint side can run on constrained edge-style hardware.

\begin{figure}[H]
  \centering
  \includegraphics[width=\textwidth]{Figures/final-architecture-report.png}
  \caption{Frozen Tier 1.5 proof-of-concept architecture.}
  \label{fig:implementation-final-architecture}
\end{figure}

\subsection{Local n8n Runtime with Docker}\label{subsec:implementation-n8n-docker}

The workflow engine was implemented as a local self-hosted n8n runtime on the PC. The runtime is defined in \codepath{infrastructure/docker/docker-compose.yml} and uses the pinned container image \texttt{n8nio/n8n:1.123.37}. The compose file exposes n8n on port \texttt{5678}, persists n8n data in a Docker volume, and uses local development settings such as disabled secure cookies and disabled diagnostics/version notifications \cite{dockerDocs,n8nDocs}.

The PC-hosted Docker setup was chosen for stability and repeatability. It allowed the workflow engine, middleware, contracts, and evaluation harness to be developed in a controlled local environment before adding Raspberry Pi validation. This also reduced the risk of turning the thesis into a deployment study of n8n on constrained hardware. Full n8n deployment on the Raspberry Pi was intentionally outside the implemented scope. The Raspberry Pi validation therefore appears later in the implementation as a split validation of the middleware/action endpoint side, not as a replacement for the PC-hosted n8n baseline.

The local n8n runtime is the basis for both implemented workflow modes. The deterministic baseline and the minimal agent-enhanced workflow both receive temperature events through n8n webhook flows and route their resulting fan action to the middleware. Keeping the workflow engine constant makes the two workflow modes easier to compare in the evaluation chapter.

\subsection{Python Middleware API Bridge}\label{subsec:implementation-middleware}

The Python middleware implements the boundary between workflow decisions and action execution. It is placed under \texttt{middleware/} and can be started with \texttt{python -m middleware.api.app}. The HTTP request handling is implemented in \codepath{middleware/api/routes.py} using Python standard-library HTTP server components \cite{pythonHttpServerDocs}, while the simulated sensor and fan state are implemented in \codepath{middleware/gpio/mock\_sensor.py}.

The middleware exposes four implemented endpoints. \texttt{GET /status} returns the middleware status and the current simulated hardware state. \texttt{POST /sensor-event} accepts a temperature event, normalizes it with a minimal inline check, stores it as the latest sensor event, and can forward it to n8n when a webhook URL is configured. \texttt{POST /fan/on} changes the simulated fan state to \texttt{on}, and \texttt{POST /fan/off} changes the simulated fan state to \texttt{off}. The fan endpoints return JSON responses that include the action name, the resulting fan state, a \texttt{simulated} flag, and a reason field.

This middleware boundary is important because it prevents workflow nodes from being treated as direct hardware controllers. Even in the Raspberry Pi validation, the action remains a simulated middleware action. No GPIO library, real fan driver, or physical actuator path is implemented in the frozen system. This keeps the implementation aligned with the thesis scope: the work evaluates an IoT-style action pipeline and contract boundary, not physical electronics control.

\subsection{Shared JSON Contracts}\label{subsec:implementation-contracts}

The shared interface layer is implemented under \codepath{shared\_interfaces/}. It defines the data structures that connect the middleware side and the workflow/agent side. The main reason for using structured JSON contracts was to avoid loose text as an action boundary. A free-form agent response or informal workflow payload would make routing and validation ambiguous. A small schema-based contract, by contrast, makes the expected fields explicit and gives the safety design a precise object to check \cite{jsonSchemaDocs}.

The input-side schema, \codepath{shared\_interfaces/json-schema/sensor-event.schema.json}, defines the temperature sensor event used by the thesis scenario. It requires \texttt{sensor\_id}, \texttt{timestamp}, \texttt{type}, \texttt{value}, and \texttt{unit}. The supported sensor type is \texttt{temperature}, and the supported unit is \texttt{C}. This intentionally small input contract keeps the scenario measurable and avoids expanding into multiple sensors or device families.

The output-side schema, \codepath{shared\_interfaces/json-schema/agent-action.schema.json}, defines the fan action object that workflow-produced actions should follow. It requires \texttt{action\_id}, \texttt{target}, \texttt{reason}, and \texttt{requires\_approval}. The allowed action identifiers are \texttt{fan\_on} and \texttt{fan\_off}, and the allowed target is \texttt{fan\_1}. The schema also disallows additional properties. These restrictions support predictable routing: after validation, \texttt{fan\_on} maps to \texttt{POST /fan/on}, and \texttt{fan\_off} maps to \texttt{POST /fan/off}.

The repository also includes example payloads for a valid sensor event, a valid fan-on action, and a blocked action example. These examples support reportability and later evaluation, but they should not be interpreted as proof of production runtime enforcement. The contract layer defines the expected boundary; the safety layer defines how the boundary should be checked.

\subsection{Deterministic Baseline Workflow}\label{subsec:implementation-deterministic}

The deterministic baseline workflow is implemented under \codepath{cognitive\_logic/workflows/} as \texttt{deterministic-baseline.md} and \texttt{deterministic-baseline.json}. It is the non-AI comparison anchor for the thesis. The workflow receives a temperature event through an n8n Webhook node, reads the event value from the parsed request body, applies one threshold rule, and then calls one of the simulated fan endpoints.

The rule is intentionally simple:

\begin{quote}
if temperature \(\geq 30.0\) C, request \texttt{fan\_on}; otherwise, request \texttt{fan\_off}.
\end{quote}

When the temperature is at or above the threshold, the workflow creates an action object corresponding to \texttt{fan\_on} and calls \texttt{POST} \codepath{http://host.docker.internal:8000/fan/on}. When the temperature is below the threshold, it creates an action object corresponding to \texttt{fan\_off} and calls \texttt{POST} \codepath{http://host.docker.internal:8000/fan/off}. The \texttt{host.docker.internal} address is used because n8n runs inside Docker while the middleware runs on the host during local PC execution.

The deterministic workflow exists for two reasons. First, it provides a predictable reference behavior for the same temperature-to-fan task used by the agent-enhanced workflow. Second, it gives the evaluation a simple baseline against which action correctness, consistency, and latency can be compared. The baseline does not use prompts, memory, model calls, safety approval nodes, GPIO, or Raspberry Pi deployment.

\begin{figure}[H]
  \centering
  \includegraphics[width=\textwidth]{Figures/deterministic-baseline-canvas.png}
  \caption{Deterministic baseline n8n workflow canvas.}
  \label{fig:deterministic-workflow-canvas}
\end{figure}

\subsection{Minimal Agent-Enhanced Workflow}\label{subsec:implementation-agent}

The minimal agent-enhanced workflow is implemented under \codepath{cognitive\_logic/workflows/} as \texttt{agent-minimal.md} and \texttt{agent-minimal.json}. It uses the same temperature-event scenario and the same middleware fan endpoints as the deterministic baseline, but it replaces the fixed threshold decision node with a constrained LLM decision step.

The workflow receives the temperature event through a Webhook node, prepares the current event as agent input, sends it to one LLM decision step, parses the returned action object, and routes supported fan actions to the middleware. The Step 7 evidence identifies the configured model node as Google Gemini Chat Model. The workflow uses one prompt set, stored in \codepath{cognitive\_logic/prompts/system-prompt-v1.md}, and one memory choice, stored in \codepath{cognitive\_logic/memory/memory-choice-v1.md}. The chosen memory strategy is stateless execution with no memory node connected.

The prompt constrains the LLM to return a structured JSON action object rather than a free-form explanation. The expected output follows the same action contract used elsewhere in the thesis: \texttt{action\_id}, \texttt{target}, \texttt{reason}, and \texttt{requires\_approval}. The workflow then parses this structured output and routes \texttt{fan\_on} to \texttt{POST /fan/on} or \texttt{fan\_off} to \texttt{POST /fan/off}. A minimal unrouted branch exists for output outside the expected action values, but this branch is not the formal safety layer.

The purpose of this workflow is narrow. It tests whether an agent-style decision step can produce contract-valid structured action output for the same small temperature-to-fan task as the deterministic baseline. It does not implement broad autonomy, a ReAct loop \cite{reactYao2023}, a multi-agent system, model benchmarking, or persistent memory.

\begin{figure}[H]
  \centering
  \includegraphics[width=\textwidth]{Figures/agent-workflow-canvas.png}
  \caption{Minimal agent-enhanced n8n workflow canvas.}
  \label{fig:agent-workflow-canvas}
\end{figure}

\subsection{Minimum Safety and Validation Design}\label{subsec:implementation-safety}

The safety layer is implemented as a minimum design and specification layer under \codepath{safety\_layer/}. It consists of an action policy, an output validation design, a human-in-the-loop approval design, and three documented example cases. The action policy is stored in \codepath{safety\_layer/policies/action-policy-v1.md}; the parser and validation design is stored in \codepath{safety\_layer/parsers/output-validation-v1.md}; and the approval design is stored in \codepath{safety\_layer/approvals/hitl-v1.md} \cite{humanInTheLoopAmershi2014}.

The safety design is based on the shared action contract. An action is allowed to proceed without human approval only when it parses as JSON, matches the action schema, uses \texttt{fan\_on} or \texttt{fan\_off}, targets \texttt{fan\_1}, contains a non-empty reason, has no unexpected fields, and sets \texttt{requires\_approval} to \texttt{false}. A malformed action, an action with missing fields, an unknown action, an unknown target, or an action with unexpected fields is blocked by the policy. A risky or unknown action may be shown for human review as evidence, but policy version 1 does not allow an unknown action or target to be approved into a middleware call.

The repository includes one allowed case, one blocked malformed case, and one risky/unknown case. The allowed case describes a valid \texttt{fan\_on} action targeting \texttt{fan\_1}. The blocked case describes an incomplete action missing required fields. The risky case describes an invented \texttt{open\_window} action targeting an unknown device, which is outside the thesis scenario and must not reach the fan endpoints.

This safety layer matters because it defines the boundary that workflow and agent output should satisfy before action execution. However, it is not a production-grade runtime safety framework. It is also not evidence that all unsafe actions are blocked by an implemented runtime gate. In the frozen implementation, it should be interpreted as the minimum validation and approval design required to make the contract boundary explicit and discussable in the report.

\subsection{Evaluation Harness}\label{subsec:implementation-evaluation-harness}

The evaluation harness is implemented under \texttt{evaluation/} and \texttt{scripts/}. Its purpose is to make the proof-of-concept measurable without adding a heavy observability stack. The harness uses standard-library Python scripts and CSV files so that test execution and aggregation remain lightweight and transparent \cite{pythonCsvDocs,csvRfc4180}.

The test cases are defined in \codepath{evaluation/datasets/test-cases.json}. They include deterministic workflow cases, agent workflow cases, and safety-mode cases. The metric fields are documented in \codepath{evaluation/metrics/metrics-schema.md}. Raw result files are written under \codepath{evaluation/results/raw/}, while processed summaries are written under \codepath{evaluation/results/processed/}.

The collection script, \codepath{scripts/collect\_metrics.py}, posts the selected test payloads to a supplied n8n webhook URL. For workflow modes, it records the mode, test case identifier, run index, input temperature, expected action, observed action, HTTP status, success flag, latency, observed response, and error/notes fields. It also attempts to collect CPU and RAM values through \texttt{docker stats --no-stream} when Docker statistics are available \cite{dockerStatsDocs}, and it records thermal data only when a supported thermal source exists. The aggregation script, \codepath{scripts/aggregate\_results.py}, reads raw CSV files and writes processed summaries for latency, resources, safety outcomes, and baseline-versus-agent comparison.

Chapter~\ref{sec:results} reports the numerical outcomes from these CSV files. This implementation chapter therefore describes the harness structure and the measured fields, but it does not present final latency, resource, or correctness values.

\subsection{Raspberry Pi Tier 1.5 Validation Setup}\label{subsec:implementation-pi}

The Raspberry Pi work was implemented as Tier 1.5 validation rather than full deployment. In this setup, n8n stayed on the PC and continued to run inside Docker. The Raspberry Pi ran only the Python middleware/action endpoint side. The middleware was started on the Pi with the host bound to \texttt{0.0.0.0}, allowing the PC-hosted n8n workflow to call the Pi over the network using HTTP \cite{raspberryPiDocs}.

The validation reused the same temperature-to-fan scenario. High-temperature and low-temperature cases were sent through the PC-hosted n8n workflow, and the workflow called the Raspberry Pi middleware endpoints for \texttt{/fan/on} and \texttt{/fan/off}. The Pi validation evidence is stored under \codepath{evaluation/results/pi-validation/}. It includes the final workflow-run CSV, notes, startup/IP transcript, canvas and output screenshots, direct endpoint latency evidence, CPU/RAM evidence, and thermal evidence.

This split validation strengthens the edge-computing angle of the thesis without changing the core architecture. It shows that the action endpoint side of the system can be moved from the PC to the Raspberry Pi while preserving the PC-hosted workflow engine. At the same time, its limitations are important: n8n was not deployed on the Raspberry Pi, the whole system was not moved to the Pi, the fan action remained simulated, and no real GPIO hardware was controlled.

\begin{figure}[H]
  \centering
  \includegraphics[width=\textwidth]{Figures/pi-validation-fan-on-output.png}
  \caption{Raspberry Pi Tier 1.5 validation output for the high-temperature fan-on case.}
  \label{fig:implementation-pi-validation-fan-on-output}
\end{figure}

\subsection{Summary of Implementation}\label{subsec:implementation-summary}

The implemented system forms one complete proof-of-concept pipeline. A temperature event can be sent into a PC-hosted n8n workflow, evaluated either by a deterministic threshold path or by a minimal agent-enhanced path, represented through a shared JSON action contract, bounded by a minimum validation and approval design, executed through Python middleware as a simulated fan action, and recorded through a CSV-based evaluation harness. Raspberry Pi validation extends this pipeline only by moving the middleware/action endpoint to the Pi while keeping n8n on the PC. This implementation is therefore narrow by design, but it is complete enough to support the thesis evaluation of action correctness, contract-valid agent output, latency, resource behavior, and deployment limitations.
